\chapter{Zero-Knowledge Proofs}

There is a kind of magic trick that sounds impossible the first time you hear
it stated plainly: \emph{I can convince you that I know a secret, and at the end
of the conversation you will be certain that I know it, yet you will have
learned nothing whatsoever about the secret itself.} Not a hint. Not a single
bit. If you tried afterwards to convince a third party of what you saw, you
would fail, because the transcript of our conversation could have been
manufactured by anyone, knowledge or no knowledge.

That sentence packs three separate promises, and the whole chapter is really an
unpacking of them. The proof must \emph{work} when the prover is telling the
truth (completeness). It must \emph{fail} when the prover is lying (soundness).
And it must \emph{leak nothing} beyond the bare fact that the statement is true
(zero-knowledge). Each of these is a precise mathematical condition, and the
genius of the field, born in a 1985 paper by Goldwasser, Micali, and Rackoff,
was figuring out how to state the third one rigorously at all.

We will build the idea from a cartoon (a cave with a magic door), tighten it
into definitions, then construct a real protocol you could ship: Schnorr's
identification scheme, the workhorse behind a great deal of modern
authentication. We turn it non-interactive with the Fiat-Shamir transform,
which is how zero-knowledge becomes a \emph{signature}. Commitments show up as
the universal glue. And we close with a glimpse of zk-SNARKs, the technology
that lets a blockchain verify a computation it never witnessed.

\chapterline

\section{The cave: intuition before formalism}
\label{ch31:cave}

The standard parable is Jean-Jacques Quisquater's ``Ali Baba cave,'' and it is
worth telling carefully because every later definition is a translation of it.

Picture a ring-shaped cave with a single entrance. The ring splits into a left
passage and a right passage, and deep inside, where the two passages meet,
there is a magic door locked by a secret word. Peggy (the prover) claims she
knows the word. Victor (the verifier) wants to be convinced, but Peggy refuses
to just say the word out loud --- that would give it away.

So they play a game. Victor waits outside, out of sight of the fork. Peggy walks
in and picks one passage at random, left or right. Now Victor steps to the
mouth of the cave and shouts the name of a passage he wants her to come out of
--- say, ``Left!'' If Peggy went down the left passage, she just walks back out;
she did not need the door. But if she went down the right passage, the only way
to emerge on the left is to open the magic door and pass through. If she knows
the word, she can always satisfy Victor's request, whichever side he names.

\begin{intuitionbox}
Here is the hinge of the whole subject. Suppose Peggy is a fraud who does
\emph{not} know the word. She can still guess which side Victor will call and
pre-position herself there. If she guesses right, she walks out the requested
side and looks legitimate. But she only guesses right half the time. So one
round catches a liar with probability $1/2$. Run the game $40$ times with fresh
random choices, and a liar survives all $40$ rounds with probability
$2^{-40} \approx 9 \times 10^{-13}$ --- less likely than picking a specific
second out of the next thirty thousand years. Repetition turns ``maybe she got
lucky'' into ``she must know the word.''
\end{intuitionbox}

Now watch the zero-knowledge part. What did Victor actually \emph{see}? A
sequence of moments where he shouted a side and Peggy emerged from that side.
That is all. He never saw the word, never saw which passage she originally
entered (he was outside), and crucially he could have \emph{filmed a fake video
that looks identical} without Peggy knowing anything: just have an actor enter a
random side each round, and only keep the takes where the actor happened to come
out the side you planned to shout. Edit those together and you have a
convincing-looking tape of a successful proof, produced with zero knowledge of
the secret. Because a knowledgeable Peggy and a lucky actor produce
statistically identical transcripts, the transcript carries \emph{no
information} about the word. That ``I could have faked the whole thing'' move is
the beating heart of the formal definition, and it has a name: the
\emph{simulator}.

\begin{focusbox}
Three properties, one sentence each, in cave terms:
\begin{itemize}
  \item \textbf{Completeness}: if Peggy knows the word, she passes every round.
  \item \textbf{Soundness}: if she does not, each round trips her with
  probability $\ge 1/2$, so many rounds expose her.
  \item \textbf{Zero-knowledge}: a transcript of a successful run could have
  been fabricated by someone with no knowledge, so it reveals nothing.
\end{itemize}
\end{focusbox}

\begin{historybox}
Goldwasser, Micali, and Rackoff introduced interactive proofs and
zero-knowledge in ``The Knowledge Complexity of Interactive Proof-Systems''
(STOC 1985), a paper rejected several times before publication --- reviewers
struggled with the idea that a proof could convey \emph{less} than the fact it
established. The trio later shared the 2012 ACM Turing Award, in part for this
work. Quisquater and a long list of co-authors wrote the cave story in ``How to
Explain Zero-Knowledge Protocols to Your Children'' (CRYPTO 1989), which remains
the gentlest on-ramp ever written for a cryptographic primitive.
\end{historybox}

\section{Interactive proofs, made precise}
\label{ch31:interactive}

Let us replace the cave with mathematics. We fix a language $L$: a set of true
statements. For example, $L$ could be ``all graphs that have a Hamiltonian
cycle,'' or ``all numbers $y$ that are squares modulo $n$,'' or in the cases we
care about most, ``all public keys for which the prover knows the matching
private key.'' A statement $x$ is true precisely when $x \in L$.

An \emph{interactive proof system} for $L$ is a pair of algorithms, the prover
$P$ and the verifier $V$, who exchange messages. $V$ runs in probabilistic
polynomial time --- it is meant to be cheap and realistic. $P$ may be
computationally unbounded in the classical theory, though in the protocols we
build $P$ is also efficient \emph{provided it holds a secret witness} $w$. At
the end $V$ outputs accept or reject.

\begin{definition}[Interactive proof, informal]
\label{ch31:def-ip}
$(P,V)$ is an interactive proof for $L$ if:
\begin{itemize}
  \item \textbf{Completeness:} for every $x \in L$, the prover makes $V$ accept
  with probability at least $1 - \negl$ (often exactly $1$).
  \item \textbf{Soundness:} for every $x \notin L$ and \emph{every} cheating
  prover strategy $P^*$ (however clever), $V$ accepts with probability at most
  some small $s$, the \emph{soundness error}. Repetition drives $s$ toward $0$.
\end{itemize}
\end{definition}

The asymmetry is deliberate. Completeness protects the truthful prover from a
false rejection; soundness protects the verifier from \emph{every conceivable}
lying prover, including ones we have not imagined. A protocol that only stops
the liars we thought of is worthless. This is the same ``quantify over all
adversaries'' stance you have seen throughout the book.

\begin{walkbox}
Why allow $V$ to be \emph{probabilistic}? Because a deterministic verifier in a
single-prover setting buys you nothing beyond an ordinary written proof (the
class \textsf{NP}). Randomness --- the unpredictable challenge ``Left or
right?'' --- is exactly what stops the prover from preparing one canned answer in
advance. The interaction plus the coins is where the extra power lives; the
complexity class of all such interactive proofs, \textsf{IP}, turns out to equal
\textsf{PSPACE}, a celebrated 1990 result of Shamir. For cryptography we never
need that much muscle, but it tells you the framework is genuinely more
expressive than static proofs.
\end{walkbox}

\subsection*{Proof of knowledge: a stronger flavor}

For authentication we usually want something sharper than ``the statement is
true.'' We want ``the prover actually \emph{knows} a witness.'' Knowing your
private key is the whole point of logging in. This is captured by a
\emph{proof of knowledge}, and the tool that defines it is the
\emph{knowledge extractor}.

\begin{definition}[Proof of knowledge, informal]
\label{ch31:def-pok}
A protocol is a proof of knowledge for a relation $R = \set{(x,w)}$ if there is
an efficient \emph{extractor} $E$ with the following power: given the ability to
run any successful prover $P^*$ and to \emph{rewind} it --- replay it from a
chosen point with a different challenge --- $E$ outputs a valid witness $w$ with
$(x,w) \in R$. Informally: anything that can convince $V$ must possess $w$,
because we could squeeze $w$ out of it by replaying.
\end{definition}

Rewinding is a thought experiment, not something a real verifier does in a live
session. Its purpose is to argue: \emph{the only way to answer two different
challenges about the same first message is to know the secret.} We will see this
concretely in Schnorr's protocol, where two accepting transcripts that share a
first message let you solve for the private key with a single subtraction.

\section{Commitments: the cryptographic envelope}
\label{ch31:commit}

Before the headline protocol we need a building block that appears in almost
every zero-knowledge construction: the \emph{commitment scheme}. Think of it as
a sealed, tamper-evident envelope.

You \textbf{commit} to a value $m$ by computing $c = \mathsf{Commit}(m, r)$ with
fresh randomness $r$, and you hand $c$ to the other party. Later you
\textbf{open} it by revealing $(m, r)$, and they check $c = \mathsf{Commit}(m,r)$.
Two properties make the envelope useful:

\begin{itemize}
  \item \textbf{Hiding:} $c$ reveals nothing about $m$. The sealed envelope is
  opaque.
  \item \textbf{Binding:} you cannot find a second $(m', r') \ne (m, r)$ with
  the same $c$. Once sealed, the contents are fixed; you cannot swap the letter.
\end{itemize}

\begin{intuitionbox}
Commitments let you say ``I have decided on my answer'' \emph{before} you learn
the other party's challenge, without revealing the answer yet. That ordering ---
lock in your move, then see the challenge, then reveal --- is what stops a cheater
from adapting their answer to the challenge after the fact. The cave achieves
the same ordering physically: Peggy commits to a passage by walking down it
before Victor shouts.
\end{intuitionbox}

A clean number-theoretic commitment is \emph{Pedersen's}. Fix a cyclic group
$\G$ of prime order $q$ with two generators $g$ and $h$ whose discrete-log
relation nobody knows. To commit to $m \in \Z_q$, pick $r \getsr \Z_q$ and set
\[
  c \;=\; g^{m}\, h^{r}.
\]
Hiding is \emph{perfect}: because $h^r$ is uniform over the group, $g^m h^r$ is
uniform too, independent of $m$. Binding rests on the discrete-log assumption:
if you could open $c$ as both $m$ and $m'$, you would have
$g^{m}h^{r} = g^{m'}h^{r'}$, hence $h^{\,r-r'} = g^{\,m'-m}$, which hands you
$\log_g h = (m'-m)(r-r')^{-1} \bmod q$ --- the very discrete log we assumed is
hard. A simpler, hash-based commitment is $c = \Hash(m \concat r)$, binding
because $\Hash$ is collision-resistant and hiding because $\Hash$ behaves like a
random function for unknown $r$.

\begin{trapbox}
You cannot have \emph{both} perfect hiding and perfect binding at once. If $c$
truly carries zero information about $m$ (perfect hiding), then in principle
\emph{some} $(m', r')$ also maps to $c$ for every $m'$, so binding can only be
computational. Pedersen is perfectly hiding and computationally binding;
hash-based commitments are typically computationally hiding and (heuristically)
binding. Pick the side you cannot afford to lose. For a coin flip you want
binding; for a sealed-bid auction you may prize hiding.
\end{trapbox}

\begin{connectionbox}
Commitments are close cousins of hash functions (the hash-functions chapter) and
of the discrete-log assumption (the Diffie-Hellman and ElGamal chapters). The
companion algorithms book covers the cost of modular exponentiation by
square-and-multiply, which is what makes $g^m h^r$ cheap to compute; the
probability book's birthday bound is what governs how collision-resistance ---
hence binding --- degrades as the output shrinks.
\end{connectionbox}

\section{Schnorr's identification protocol}
\label{ch31:schnorr}

Now the centerpiece. Schnorr's scheme lets Peggy prove she knows the discrete
logarithm of her public key without revealing it. It is short, elegant, and
underlies real signature standards (EdDSA among them). We work in a cyclic group
$\G = \langle g \rangle$ of prime order $q$. Peggy's secret is $x \getsr \Z_q$;
her public key is $y = g^x$. The statement she proves: ``I know $x$ such that
$g^x = y$.''

The protocol is a three-move pattern cryptographers call a \emph{Sigma
protocol}, after the shape of the message flow: commitment, challenge, response.

\begin{enumerate}
  \item \textbf{Commitment.} Peggy picks a random nonce $k \getsr \Z_q$ and sends
  $a = g^{k}$. (She commits to a fresh ephemeral exponent.)
  \item \textbf{Challenge.} Victor picks $e \getsr \Z_q$ (or from a smaller
  challenge space) and sends $e$.
  \item \textbf{Response.} Peggy computes $s = k + e\,x \bmod q$ and sends $s$.
  \item \textbf{Verify.} Victor accepts iff $g^{s} = a \cdot y^{e}$.
\end{enumerate}

\begin{center}
\begin{tikzpicture}[font=\small,>=Latex]
  \node[draw,rounded corners,fill=aliceblue!12,minimum width=3.6cm,minimum height=0.9cm] (P) at (0,0) {Peggy (knows $x$, $y=g^x$)};
  \node[draw,rounded corners,fill=bobgreen!12,minimum width=3.4cm,minimum height=0.9cm] (V) at (9,0) {Victor (knows $y$)};
  \draw[->,thick] (P.south) ++(0,-0.2) -- ++(0,-3.0);
  \draw[->,thick] (V.south) ++(0,-0.2) -- ++(0,-3.0);
  \node[anchor=west] at ($(P.south)+(0,-0.55)$) {pick $k\getsr\Z_q$};
  \draw[->,thick] ($(P.south)+(0,-1.0)$) -- node[above]{$a = g^{k}$} ($(V.south)+(0,-1.0)$);
  \node[anchor=east] at ($(V.south)+(0,-1.45)$) {pick $e\getsr\Z_q$};
  \draw[<-,thick] ($(P.south)+(0,-1.9)$) -- node[above]{$e$} ($(V.south)+(0,-1.9)$);
  \node[anchor=west] at ($(P.south)+(0,-2.35)$) {$s = k+ex \bmod q$};
  \draw[->,thick] ($(P.south)+(0,-2.85)$) -- node[above]{$s$} ($(V.south)+(0,-2.85)$);
  \node[anchor=east,align=right] at ($(V.south)+(0,-3.35)$) {accept iff $g^{s}=a\,y^{e}$};
\end{tikzpicture}
\end{center}

\subsection*{Why verification works (completeness)}

If Peggy follows the protocol, the check is just algebra:
\[
  g^{s} \;=\; g^{\,k + e x} \;=\; g^{k}\,(g^{x})^{e} \;=\; a\cdot y^{e}.
\]
So a rule-following Peggy always passes. Completeness holds with probability $1$.

\subsection*{A worked example with tiny numbers}
\label{ch31:schnorr-example}

Let us trace it concretely in a group small enough to do by hand. Take the
multiplicative group modulo the prime $p = 23$. Its order is $22 = 2 \cdot 11$.
The element $g = 2$ has order $11$ (you can check $2^{11} \equiv 1 \pmod{23}$ and
no smaller power is $1$), so $g$ generates a subgroup $\G$ of prime order
$q = 11$. We do all exponent arithmetic mod $q = 11$ and all group arithmetic
mod $p = 23$.

Peggy's secret: $x = 7$. Her public key: $y = g^{x} = 2^{7} = 128 \equiv 13
\pmod{23}$. So $y = 13$.

\begin{enumerate}
  \item Peggy draws nonce $k = 4$. She sends $a = 2^{4} = 16 \pmod{23}$.
  \item Victor draws challenge $e = 6$.
  \item Peggy responds $s = k + e x = 4 + 6\cdot 7 = 46 \equiv 2 \pmod{11}$.
        She sends $s = 2$.
  \item Victor checks. Left side: $g^{s} = 2^{2} = 4$. Right side:
        $a\cdot y^{e} = 16 \cdot 13^{6} \pmod{23}$. Compute $13^{2}=169\equiv 8$,
        $13^{3}=13\cdot 8 = 104 \equiv 12$, $13^{6}=(13^{3})^{2}=12^{2}=144
        \equiv 6 \pmod{23}$. So the right side is $16\cdot 6 = 96 \equiv 4
        \pmod{23}$.
\end{enumerate}
Both sides equal $4$. Victor accepts, and at no point did $x = 7$ cross the
wire; he only saw $(a,e,s) = (16, 6, 2)$.

\begin{labbox}
Re-run the example yourself with $k = 9$ and $e = 3$. You should get
$s = 9 + 21 = 30 \equiv 8 \pmod{11}$, $a = 2^{9} = 512 \equiv 6 \pmod{23}$, and
both sides of the check equal to $2^{8} = 256 \equiv 3 \pmod{23}$. Confirm the
right side: $y^{3} = 13^{3} \equiv 12$, times $a = 6$ gives $72 \equiv 3
\pmod{23}$. It matches --- a different transcript, same accept.
\end{labbox}

\subsection*{Why a cheat fails (soundness via extraction)}

Here is the elegant argument. Suppose a prover, knowing only $y$ but not $x$,
can answer \emph{two different challenges} $e \ne e'$ for the \emph{same}
first message $a$, producing valid responses $s$ and $s'$. Both pass
verification:
\[
  g^{s} = a\, y^{e}, \qquad g^{s'} = a\, y^{e'}.
\]
Divide one by the other: $g^{\,s - s'} = y^{\,e - e'}$. Since $y = g^{x}$, this
says $g^{\,s-s'} = g^{\,x(e-e')}$, so in the exponent (mod $q$):
\[
  s - s' \equiv x\,(e - e') \pmod{q}
  \quad\Longrightarrow\quad
  x \;=\; (s - s')\,(e - e')^{-1} \bmod q.
\]
The inverse exists because $q$ is prime and $e \ne e'$. \textbf{Anyone who can
answer two challenges on one commitment can compute $x$ outright.} That is the
knowledge extractor of Definition~\ref{ch31:def-pok} made fully concrete: rewind
the prover to just after it sent $a$, feed a fresh challenge, and subtract.

\begin{whybox}
\textbf{This is why Schnorr is a proof of knowledge.} A prover who does not know
$x$ cannot consistently answer two challenges on the same $a$ --- if it could, the
extraction above would compute $x$ from the prover's own outputs, contradicting
the hardness of discrete log. So with challenge space of size $q$, the best a
cheater can do is guess the single challenge it can answer and pray Victor sends
exactly that one: soundness error $1/q$. With a $128$-bit prime $q$, that is
$2^{-128}$ in one shot. (With a tiny challenge set, say a single bit
$e\in\set{0,1}$, the error is $1/2$ per round, and you repeat --- exactly the
cave.)
\end{whybox}

\subsection*{Why it leaks nothing (zero-knowledge via simulation)}

Completeness and soundness are about correctness. Zero-knowledge is the
delicate one, and we prove it by building the simulator from the cave: a machine
that produces transcripts \emph{indistinguishable} from real ones \emph{without
knowing $x$}.

For \emph{semi-trusted-verifier} zero-knowledge (the verifier picks $e$ at
random, exactly as the protocol specifies, rather than adversarially) the
simulator is startlingly simple. To fake a transcript $(a,e,s)$:

\begin{enumerate}
  \item Pick the response $s \getsr \Z_q$ first.
  \item Pick the challenge $e \getsr \Z_q$.
  \item \emph{Solve for} the commitment that makes the check pass:
        $a = g^{s}\, y^{-e}$.
\end{enumerate}

By construction $g^{s} = a\,y^{e}$, so the fake transcript verifies. And its
distribution is identical to a real run: in a real run $k$ is uniform so $a$ is
uniform and $s = k + ex$ is uniform; in the simulation $s$ is uniform and $e$ is
uniform, which forces $a$ to be uniform too. Same distribution, no secret used.

\begin{intuitionbox}
The simulator cheats by working the protocol \emph{out of order}. The real
prover commits, then is challenged, then must respond --- and to respond it needs
$x$. The simulator, free to choose $e$ and $s$ itself, runs time backwards:
decide the answer and the challenge, then retrofit a commitment that fits. A
real verifier cannot do this because it does not get to choose $a$ after seeing
$e$ and $s$. But the resulting \emph{transcript} looks identical, which is the
whole point: if a perfect forgery exists, the transcript cannot contain
information about $x$.
\end{intuitionbox}

\begin{formalbox}
\textbf{Zero-knowledge against a rule-following verifier (the standard
literature name is HVZK).} A Sigma protocol has this property if there is a
PPT simulator $S$ that, on input the statement $y$ (no witness), outputs
transcripts $(a,e,s)$ whose distribution equals (or is computationally
indistinguishable from) that of real protocol runs in which the verifier draws
$e$ uniformly as specified. For Schnorr the two distributions are
\emph{identical}, so the simulation is \emph{perfect}. Upgrading to security
against a \emph{malicious} verifier who picks $e$ adversarially needs more
care --- one route is to have the verifier first commit to $e$ before seeing
$a$; another is the Fiat-Shamir transform of the next section, which removes the
verifier's freedom entirely by deriving $e$ from a hash.
\end{formalbox}

\begin{trapbox}
\textbf{Never reuse the nonce $k$.} If Peggy ever signs or proves twice with the
same $a = g^k$ under two different challenges $e_1 \ne e_2$, an eavesdropper runs
exactly the extractor above on the two public transcripts and recovers $x$:
$x = (s_1 - s_2)(e_1 - e_2)^{-1}$. This is not hypothetical --- it is the same
flaw that leaked the Sony PlayStation~3 ECDSA signing key in 2010 (a fixed nonce)
and has burned countless wallets. The nonce must be fresh, uniform, and secret
every single time. Deterministic nonce generation (RFC 6979, or hashing the
message with the key) is the modern defense.
\end{trapbox}

\section{Fiat-Shamir: from conversation to signature}
\label{ch31:fiat-shamir}

Interaction is inconvenient. A signature, a login token, a blockchain proof ---
these should be one-shot artifacts you can post and walk away from, with no live
verifier to volley challenges. The \emph{Fiat-Shamir transform} (Amos Fiat and
Adi Shamir, 1986) removes the verifier from the loop with a single brilliant
substitution:

\begin{center}
\emph{Replace the verifier's random challenge $e$ with}
$e = \Hash(\text{everything so far})$.
\end{center}

The prover can compute the hash herself, so she no longer needs anyone to talk
to. But because a good hash is unpredictable, she still cannot choose $a$
\emph{after} seeing $e$ --- the hash of $a$ \emph{is} $e$, so fixing $a$ fixes
$e$. The unpredictability that the live verifier used to supply is now supplied
by the hash function. To prevent the soundness argument from collapsing, we
model $\Hash$ as a random oracle.

For Schnorr, the transformed, non-interactive protocol on statement $y = g^x$ is:

\begin{enumerate}
  \item Pick $k \getsr \Z_q$, set $a = g^{k}$.
  \item Compute $e = \Hash(g \concat y \concat a)$. \emph{(For a signature on
        message $M$, hash $M$ in too:} $e = \Hash(g \concat y \concat a \concat
        M)$.\emph{)}
  \item Compute $s = k + e x \bmod q$.
  \item Output the proof or signature $\sigma = (a, s)$ (equivalently $(e,s)$).
\end{enumerate}
A verifier recomputes $e = \Hash(g \concat y \concat a \concat M)$ and checks
$g^{s} = a\, y^{e}$. That is the \textbf{Schnorr signature scheme} --- and EdDSA
(Ed25519), used in SSH, TLS 1.3, Signal, and many cryptocurrencies, is precisely
this idea over a specific elliptic curve.

\begin{whybox}
\textbf{Why hashing the message in gives a secure signature.} Binding $e$ to $M$
means a forger who wants a signature on a chosen $M$ is forced into the same
trap as a cheating prover: to produce a valid $(a,s)$ for that $M$ they would
need to answer the challenge $e = \Hash(\dots\concat M)$ on a commitment $a$,
which by the soundness/extraction argument requires knowing $x$. In the random
oracle model one proves this rigorously via a ``forking lemma'': run the forger,
rewind it to the point it queried the oracle on $a$, answer with a fresh $e'$,
and obtain two transcripts $(a,e,s)$, $(a,e',s')$ --- then extract $x$. A forger
thus implies a discrete-log solver.
\end{whybox}

\begin{trapbox}
\textbf{Always hash the statement, not just the commitment.} A notorious class
of bugs (``weak Fiat-Shamir'') comes from hashing too little --- e.g.\ computing
$e = \Hash(a)$ but forgetting to bind $y$ (the public key or statement) and any
public parameters. An attacker who can choose the statement \emph{after} seeing
the challenge can then forge. This exact omission produced real CVEs in
production proof systems (the 2022 ``Frozen Heart'' vulnerabilities affecting
several SNARK and Bulletproof libraries). Rule: the hash input must commit to
\emph{the full statement and all public values}, not just the prover's first
message.
\end{trapbox}

\begin{implementbox}
A self-contained Schnorr identification plus Fiat-Shamir signature over a small
prime-order subgroup. \textbf{Learning only --- do not use for anything real.}
The group is microscopic, the randomness is Python's non-cryptographic PRNG, and
there is no constant-time arithmetic or input validation.
\begin{lstlisting}[language=Python]
import hashlib, random

# Toy prime-order group: subgroup of Z_p^* of order q, generator g.
p, q, g = 2267, 103, 354      # 354 has order 103 modulo 2267 (toy params)

def H(*parts):                # Fiat-Shamir challenge in [0, q)
    data = "|".join(str(x) for x in parts).encode()
    return int.from_bytes(hashlib.sha256(data).digest(), "big") % q

# --- key generation ---
x = random.randrange(1, q)    # secret key
y = pow(g, x, p)              # public key  y = g^x mod p

# --- interactive Schnorr proof of knowledge of x ---
def prove_interactive(x):
    k = random.randrange(1, q)
    a = pow(g, k, p)                  # 1. commitment
    e = random.randrange(0, q)        # 2. verifier's challenge
    s = (k + e * x) % q               # 3. response
    return a, e, s

def verify_interactive(y, a, e, s):
    return pow(g, s, p) == (a * pow(y, e, p)) % p

a, e, s = prove_interactive(x)
assert verify_interactive(y, a, e, s)

# --- non-interactive signature on a message via Fiat-Shamir ---
def sign(x, msg):
    k = random.randrange(1, q)        # MUST be fresh & secret each time
    a = pow(g, k, p)
    e = H(g, y, a, msg)               # challenge from the hash, binds y and msg
    s = (k + e * x) % q
    return a, s

def verify(y, msg, sig):
    a, s = sig
    e = H(g, y, a, msg)
    return pow(g, s, p) == (a * pow(y, e, p)) % p

sig = sign(x, "pay Bob 5 coins")
assert verify(y, "pay Bob 5 coins", sig)
assert not verify(y, "pay Bob 500 coins", sig)   # message is bound

# --- the nonce-reuse catastrophe, demonstrated ---
def sign_fixed_nonce(x, msg, k):       # DON'T: reuse k across messages
    a = pow(g, k, p)
    e = H(g, y, a, msg)
    return a, (k + e * x) % q, e
a1, s1, e1 = sign_fixed_nonce(x, "msg one", k=42)
a2, s2, e2 = sign_fixed_nonce(x, "msg two", k=42)   # same k -> same a
# Two transcripts, one commitment: recover the secret by subtraction.
x_recovered = ((s1 - s2) * pow(e1 - e2, -1, q)) % q
assert x_recovered == x                # secret key fully leaked
print("recovered secret:", x_recovered, "== real:", x)
\end{lstlisting}
The last block is the lesson that matters most: a single repeated nonce turns a
secure scheme into a public broadcast of your private key.
\end{implementbox}

\section{A second example: graph 3-coloring}
\label{ch31:3color}

Schnorr proves an algebraic fact. But zero-knowledge is far more general --- a
classic theorem says \emph{every} statement in \textsf{NP} has a zero-knowledge
proof. The bridge is graph 3-coloring, an \textsf{NP}-complete problem, and its
protocol is the cleanest illustration of commitments doing the heavy lifting.

A proper 3-coloring of a graph paints each vertex red, green, or blue so that no
edge joins two same-colored vertices. Peggy knows such a coloring and wants to
convince Victor the graph is 3-colorable without revealing which colors go
where.

\begin{enumerate}
  \item Peggy randomly \emph{permutes} the three color names (so the actual
  colors are scrambled this round), then \textbf{commits} to the color of every
  vertex --- one sealed envelope per vertex.
  \item Victor picks one \emph{edge} at random and asks Peggy to open the two
  envelopes at its endpoints.
  \item Peggy opens those two commitments. Victor checks they are valid openings
  and that the two colors \emph{differ}. He accepts this round iff they do.
\end{enumerate}

\begin{center}
\begin{tikzpicture}[font=\small,>=Latex,
   v/.style={circle,draw,minimum size=7mm,thick}]
  \node[v,fill=red!25]   (A) at (0,1.4)   {$A$};
  \node[v,fill=green!25] (B) at (1.6,0.4) {$B$};
  \node[v,fill=blue!25]  (C) at (0.9,-1.2){$C$};
  \node[v,fill=red!25]   (D) at (-0.9,-1.2){$D$};
  \node[v,fill=blue!25]  (E) at (-1.6,0.4){$E$};
  \draw[thick] (A)--(B) (B)--(C) (C)--(D) (D)--(E) (E)--(A) (A)--(C) (B)--(D);
  \node[align=left,anchor=west] at (3.0,0.2)
    {Every edge joins two\\different colors.\\
     Victor opens one edge,\\sees two distinct colors,\\learns nothing global.};
\end{tikzpicture}
\end{center}

\begin{whybox}
\emph{Completeness:} a real coloring has differently-colored endpoints on every
edge, so Peggy always passes. \emph{Soundness:} if the graph is \emph{not}
3-colorable, then under \emph{any} commitment Peggy makes, at least one edge has
matching endpoints. Victor hits that bad edge with probability at least $1/|E|$,
so each round catches a fraud with probability $\ge 1/|E|$; repeating
$|E|\cdot 40$ rounds drives the soundness error below $2^{-40}$.
\emph{Zero-knowledge:} Victor learns only that \emph{two specific vertices have
different colors} --- which, because Peggy re-permutes the colors every round, is
just ``two random distinct colors,'' carrying no information about the global
coloring. The simulator fakes a round by committing to two random distinct
colors on the queried edge and garbage elsewhere; since commitments are hiding,
Victor cannot tell the difference.
\end{whybox}

The binding property of the commitments is what nails Peggy down: she must fix
all colors \emph{before} Victor names the edge, so she cannot recolor on the fly
to dodge a bad edge. The hiding property is what protects her secret. That
interplay --- bind to stop cheating, hide to protect knowledge --- recurs in
every interactive zero-knowledge proof.

\section{A glimpse of zk-SNARKs}
\label{ch31:snark}

Schnorr proves knowledge of one discrete log. Real applications want to prove
sprawling statements --- ``I ran this 10,000-step program correctly,'' ``this
transaction is valid under rules nobody should see my inputs to'' --- and they
want the proof to be \emph{tiny} and \emph{instant to check}, even if the
computation it attests to was enormous. That is the promise of a
\textbf{zk-SNARK}: \emph{Zero-Knowledge Succinct Non-interactive ARgument of
Knowledge.} Unpack the acronym and you have the whole spec.

\begin{itemize}
  \item \textbf{Succinct:} the proof is small (often a few hundred bytes) and
  verification is fast (milliseconds), \emph{independent of how big the
  computation was}. This is the headline feature --- you can verify in moments a
  computation that took an hour to run.
  \item \textbf{Non-interactive:} one message, no back-and-forth (Fiat-Shamir or
  a structured reference string makes this happen).
  \item \textbf{ARgument:} soundness holds against \emph{computationally
  bounded} provers only --- a small but real weakening from a ``proof,'' the
  price of succinctness.
  \item \textbf{of Knowledge:} as with Schnorr, a valid proof implies the prover
  actually holds a witness.
  \item \textbf{Zero-Knowledge:} the proof leaks nothing about the witness.
\end{itemize}

The machinery underneath is heavier than anything else in this chapter, so we
only sketch the pipeline. First the computation is written as an
\emph{arithmetic circuit} over a finite field, then flattened into a
\emph{rank-1 constraint system} (R1CS): a list of equations of the form
$(\mathbf{a}\cdot\mathbf{z})(\mathbf{b}\cdot\mathbf{z}) =
(\mathbf{c}\cdot\mathbf{z})$,
where $\mathbf{z}$ holds all the wire values including the secret witness. Those
constraints are then encoded as polynomials (a QAP, quadratic arithmetic
program), and the prover's job becomes ``demonstrate that a certain polynomial
identity holds'' --- something that can be checked at a \emph{single random
point} thanks to the fact that a nonzero low-degree polynomial has few roots.
\emph{Polynomial commitments} (KZG commitments, built on elliptic-curve
pairings, or hash-based ones in STARKs) let the prover commit to those
polynomials and reveal just one evaluation, which is where ``succinct'' comes
from.

\begin{intuitionbox}
The unifying trick is the same one that made Schnorr's soundness work, scaled
up. To prove ``these two giant polynomials are equal,'' you do \emph{not} send
the polynomials. You let the verifier pick a random point and you reveal both
values there. If the polynomials differ, they agree at only a handful of points
out of the whole field, so a random point catches the lie with overwhelming
probability (Schwartz-Zippel). Commitments stop the prover from choosing the
polynomial after seeing the point. Random challenge plus commitment plus
``few roots'' equals soundness --- the exact recipe from the cave, just with
polynomials instead of passages.
\end{intuitionbox}

\begin{trapbox}
\textbf{The trusted-setup landmine.} Many SNARKs (Groth16, the classic
zk-SNARK) require a one-time \emph{trusted setup} that produces public
parameters from secret randomness --- the ``toxic waste.'' Whoever knows that
randomness can forge proofs of false statements at will. Real deployments run
elaborate multi-party ``ceremonies'' (Zcash's ``Powers of Tau'' had thousands of
participants) so that the secret is safe as long as \emph{one} participant
deleted their share. STARKs and Bulletproofs avoid trusted setup entirely, at
the cost of larger proofs. When you evaluate a SNARK system, ask first: who held
the toxic waste, and do I trust that it was destroyed?
\end{trapbox}

\begin{applicationbox}
\textbf{Where zero-knowledge actually ships.}
\begin{itemize}
  \item \textbf{Authentication.} Schnorr-style proofs of knowledge let you log
  in by proving you hold a private key without sending it; this is the spirit of
  challenge-response in SSH public-key auth and of FIDO2/WebAuthn hardware keys.
  \item \textbf{Signatures.} EdDSA (Ed25519) is Fiat-Shamir-transformed Schnorr,
  used in TLS 1.3, OpenSSH, Signal, and the Bitcoin Taproot upgrade.
  \item \textbf{Private cryptocurrency.} Zcash uses zk-SNARKs so that a payment
  can be verified valid --- inputs equal outputs, no double-spend --- while the
  sender, receiver, and amount stay encrypted.
  \item \textbf{Blockchain scaling (zk-rollups).} Systems like zkSync,
  StarkNet, and Polygon zkEVM execute thousands of transactions off-chain and
  post a single succinct proof that all of them were valid, so the main chain
  verifies a batch in one cheap check instead of re-running every transaction.
  \item \textbf{Privacy-preserving identity.} ``Prove you are over 18'' or
  ``prove your credential is in the valid set'' without revealing your birthdate
  or identity --- the core of anonymous-credential and selective-disclosure
  systems.
\end{itemize}
\end{applicationbox}

\section{Putting the pieces together}
\label{ch31:summary}

Strip away the machinery and one shape remains. A zero-knowledge proof is a
disciplined conversation governed by three guarantees you can now state without
hand-waving. \emph{Completeness} is verified by algebra: the rule-following
prover's response always satisfies the check. \emph{Soundness} is verified by extraction:
two answers to two challenges on one commitment yield the secret, so anyone who
can pass must possess it. \emph{Zero-knowledge} is verified by simulation: a
machine with no secret produces identical-looking transcripts, so the
transcripts cannot carry the secret. The Sigma-protocol triple of
\emph{commit, challenge, respond} is the skeleton; commitments supply the
binding that stops cheating and the hiding that protects knowledge; Fiat-Shamir
freezes the live conversation into a postable artifact by replacing the
challenge with a hash; and zk-SNARKs scale the same logic --- random challenge,
commitment, few-roots soundness --- up to arbitrary computations checked in
milliseconds.

\begin{recoverybox}
Before the exercises, make sure you can answer these in your own words:
(1) In Schnorr, why does answering two challenges on the same $a$ reveal $x$?
(2) Why is reusing the nonce $k$ catastrophic, and what is the one-line attack?
(3) What does the simulator do, and why does its existence prove the protocol
leaks nothing? (4) In Fiat-Shamir, what supplies the unpredictability that the
live verifier used to provide, and why must the message be hashed in?
\end{recoverybox}

\section*{Exercises}
\addcontentsline{toc}{section}{Exercises}

\begin{exercise}[Trace Schnorr]
Work modulo $p = 23$ with $g = 2$ (order $q = 11$) and secret $x = 5$, so
$y = 2^{5} \equiv 9 \pmod{23}$. Peggy picks nonce $k = 8$ and Victor sends
challenge $e = 4$. Compute $a$, the response $s$, and verify that
$g^{s} \equiv a\,y^{4} \pmod{23}$ on both sides. State the numeric value both
sides share.
\end{exercise}

\begin{exercise}[Break a reused nonce]
A careless prover (params $g=2$, $q=11$, $p=23$) publishes two Schnorr
signatures that reuse the same nonce, hence the same $a = 6$: first
$(e_1, s_1) = (3, 2)$, second $(e_2, s_2) = (7, 0)$. Recover the secret $x$ using
$x = (s_1 - s_2)(e_1 - e_2)^{-1} \bmod 11$. Then verify your $x$ by checking that
the public key $y = g^{x} \bmod 23$ makes both transcripts pass
$g^{s} \equiv a\,y^{e} \pmod{23}$.
\end{exercise}

\begin{exercise}[Build the simulator]
For Schnorr with public key $y$, you are handed challenge $e$ in advance.
Describe, with formulas, how to produce an accepting transcript $(a,e,s)$ without
knowing $x$. Prove that your faked $a$ is uniformly distributed over $\G$, and
explain why this implies the protocol reveals nothing about $x$ to a
verifier who draws $e$ as specified.
\end{exercise}

\begin{exercise}[Soundness error and repetition]
A Sigma protocol uses a single-bit challenge $e \in \set{0,1}$, giving soundness
error $1/2$ per round. (a) How many independent rounds are needed to push the
total soundness error below $2^{-80}$? (b) If instead the challenge is drawn
uniformly from $\Z_q$ with $q$ a $256$-bit prime, what is the soundness error of
a \emph{single} round, and why is repetition then unnecessary?
\end{exercise}

\begin{exercise}[Pedersen binding implies discrete log]
Recall the Pedersen commitment $c = g^{m} h^{r}$ in a group of prime order $q$.
Suppose an adversary outputs two distinct openings $(m,r) \ne (m',r')$ of the
same $c$. Show step by step how to compute $\log_g h \bmod q$ from these
openings. Why does this prove that breaking binding is at least as hard as
discrete log?
\end{exercise}

\begin{exercise}[Why permute the colors?]
In the graph 3-coloring protocol, suppose Peggy \emph{forgets} to randomly
permute the three colors before committing each round, instead committing to the
same fixed coloring every time. Describe concretely what an adversarial Victor
learns after $O(|E|)$ rounds, and argue that the protocol is no longer
zero-knowledge.
\end{exercise}

\begin{exercise}[Weak Fiat-Shamir]
A flawed implementation computes the challenge as $e = \Hash(a)$, omitting the
statement $y$ from the hash. Sketch how an attacker who is free to choose the
public key $y$ \emph{after} computing $a$ and $e$ can produce an accepting
transcript without knowing any discrete log. (Hint: pick $s$ and $e$ first, set
$a = g^{s} y^{-e}$ --- but now you also get to pick $y$.) Explain why including
$y$ in the hash input closes the hole.
\end{exercise}

\begin{exercise}[Completeness proof]
Prove the completeness of Schnorr in full generality: for any secret $x$, nonce
$k$, and challenge $e$ in $\Z_q$, the prescribed response $s = k + ex \bmod q$
satisfies $g^{s} = a\,y^{e}$ where $a = g^{k}$ and $y = g^{x}$. State precisely
which group law you use at each step.
\end{exercise}

\begin{exercise}[Designing a challenge space]
You are deploying Fiat-Shamir Schnorr and must choose the output range of the
hash that produces $e$. Argue for the minimum bit-length of $e$ needed so that
the soundness error $1/|\text{challenge space}|$ is below $2^{-128}$. Then
explain why making $e$ \emph{larger} than $q$ (the group order) before reducing
mod $q$ does not increase security but can introduce a subtle bias --- and how to
avoid that bias.
\end{exercise}

\begin{exercise}[zk-SNARK threat model, conceptual]
A zk-rollup posts a succinct proof that a batch of $10{,}000$ transactions was
executed correctly, and the layer-1 chain verifies it in a few milliseconds.
(a) Which of the three properties --- completeness, soundness, zero-knowledge ---
is the load-bearing one for \emph{the chain's} security, and why? (b) For a
SNARK requiring trusted setup, explain in two or three sentences what the
``toxic waste'' is and the precise condition under which the system remains
sound. (c) Name one tradeoff you accept by choosing a transparent
(setup-free) STARK instead.
\end{exercise}
